\section{创新定位与综合评价}

\subsection{核心贡献链条}
项目的创新不应表述为“使用了一个神经网络”，而应集中在以下系统性结构：
\begin{equation}
\boxed{
\text{validated Maxwell physics}
\rightarrow
\text{energy-consistent }Z/D/H\text{ extraction}
\rightarrow
\text{deterministic geometry-aware thermal coordinates}
\rightarrow
\text{geometry-parametric neural EM surrogate}
\rightarrow
\text{joint feasible physical decode}
\rightarrow
\text{exact current/circuit physics}
\rightarrow
\text{true thermal ROM dynamics}.
}
\end{equation}

神经网络替代的是复杂的 geometry-to-EM-tensor map；几何对局部 thermal structures 的平移/旋转作用则显式保留为 deterministic coordinate transform。这样既避免 online Maxwell solve，也避免用一个固定全局 thermal basis 记忆所有移动热点。

\subsection{为什么学习 $Z/D_{\rm vol}/H$}
当前材料模型下 Maxwell operator 只依赖 geometry。固定 geometry 后，一次 multi-RHS solve 同时确定整个端口 current space，因此学习
\begin{equation}
g\mapsto\{Z_{\rm field},D_{\rm vol},H_j\}
\end{equation}
具有以下优势：
\begin{enumerate}
  \item target 与 impedance、total loss 和 reduced thermal source 直接相关；
  \item current magnitude/phase 由 exact quadratic forms 覆盖，无需把 current 当 neural sampling dimension；
  \item $D_{\rm vol}$ 与 $H_j$ 位于统一 Joule energy structure；
  \item geometry-dependent Loewner bounds 把 modal source 与 total dissipation 联系起来；
  \item thermal trajectory 中不重复计算高维 electromagnetic field。
\end{enumerate}

\subsection{为什么 thermal basis 必须 geometry-aware}
当 TX/RX 在固定 Eulerian grid 上发生大范围 translation/rotation 时，短时 heat source 与温度 hotspot 会跟随设备移动。若采用一个固定全局 linear basis，basis 必须同时包含“热点位于所有可能位置”的方向，导致 rank 随 geometry coverage 快速增长，并可能出现 train anchors 已被记住、held-out geometry 仍接近 span 外的现象。

因此 production theory 使用
\begin{equation}
\boxed{
\Phi(g)=
[\Phi_{\rm bg},\mathcal T_{\rm tx}(g)\Psi_{\rm tx},\mathcal T_{\rm rx}(g)\Psi_{\rm rx}].
}
\end{equation}
已知 rigid motion 由 coordinate transform 处理，而不是转化成 linear-subspace rank。该设计直接对应系统的几何物理，并比 neural basis prediction、Grassmann interpolation 或大量 global-mode enrichment 更简单。

\subsection{为什么不是 trajectory surrogate}
直接学习
\begin{equation}
(g,c,a_0,t)\mapsto T(t)
\end{equation}
需要网络同时学习 initial conditions、time scales、thermal diffusion、current scaling、circuit feedback 与 steady structure。当前方案把这些已知结构移出 learning problem，只留下静态 geometry-to-tensor map；thermal dynamics 继续由真实 $M_r(g),K_r(g)$ 控制。

\subsection{时间尺度处理的定位}
thermal ROM 不以“覆盖尽可能宽的 resolvent time range”为目标，而以“覆盖 full-order discretization 真正可解析、应用真正需要的 thermal dynamics”为目标。若 $\sqrt{\alpha\tau}$ 远小于 thermal mesh resolution，则极短时 anchor 主要在逼迫 ROM 表示离散 source placement。

production basis 优先覆盖可解析短/中时尺度和 $s=0$ steady limit；任意更长 finite-time horizon 由同一个 stable reduced ODE 连续/链式推进到达。这把“模型空间覆盖范围”和“用户查询时间长度”从理论上分离。

\subsection{结构保持的严格含义}
本文所称 structure-preserving 只指已经按构造或理论明确保证的结构：
\begin{itemize}
  \item reciprocal assumptions 下的 $Z^T=Z$；
  \item $D_{\rm vol}\succeq0$ 与 $\operatorname{Herm}Z-D_{\rm vol}\succeq0$；
  \item $H_j=H_j^H$ 与 geometry-dependent modal Loewner bounds；
  \item exact complex-current quadratic scaling；
  \item deterministic canonical mode ordering 与 geometry-aware thermal coordinate construction；
  \item 真实 $M_r(g),K_r(g)$ 与静态 geometry 零源 thermal dissipation；
  \item explicit initial-condition/circuit/steady-state semantics。
\end{itemize}
这些条件不等价于 exact spatial realizability、pointwise positivity、全参数域 surrogate certificate 或 global nonlinear stability。

\subsection{理论强项与风险}
\begin{center}
\begin{tabular}{p{0.19\linewidth}p{0.72\linewidth}}
\toprule
指标 & 评价与风险 \\
\midrule
理论分层 & 强。Maxwell、EM tensors、deterministic thermal coordinates、circuit 与 thermal dynamics 的责任边界明确。 \\
能量结构 & 强。truth 使用独立 Poynting audit；surrogate 使用联合 $Z/D/H$ feasible set。 \\
几何 thermal 泛化 & 相比固定 global basis 更合理。rigid motion 被显式 quotient 掉，不再消耗大量 linear rank。 \\
学习任务 & 中到强。网络只学习 geometry-to-EM-tensor map；可学习性仍取决于 geometry discretization continuity、POD spectrum 与 domain coverage。 \\
thermal ROM & 条件强。需要 transport/interpolation、basis conditioning、canonical rank 与 held-out full trajectory audit 全部通过。 \\
推理速度 & 强。静态 geometry 只需一次 MLP、一次 deterministic small thermal assembly，随后是 small ODE/circuit operations。 \\
可验证性 & 强于 trajectory black box。tensor、power balance、modal bounds、transport error、ROM error 与 trajectory error 可分层审计。 \\
模型风险 & source regularization、feed/return path、wire internal impedance、thermal convection、size/deformation beyond rigid transport 仍需验证。 \\
\bottomrule
\end{tabular}
\end{center}

\subsection{正式报告应回答的问题}
\begin{enumerate}
  \item 如何建立 reciprocity、passivity、independent Poynting audit 与 thermal Joule consistency 同时成立的 $Z/D/H$ objects？
  \item 如何利用 multi-RHS linearity 与 Hermitian current structure 消除 operating-current sampling dimension？
  \item rigid geometry motion 从 fixed global basis 中剥离后，canonical thermal mode rank 与 held-out error 能下降多少？
  \item geometry-aware transport 的 interpolation/conditioning error 与真正 thermal operator variation 各占多少？
  \item geometry-to-$Z/D/H$ map 是否具有足够低的有效复杂度，可被紧凑 surrogate 学习？
  \item tensor surrogate、thermal ROM、transport、initial-condition 与 circuit errors 如何传播到 trajectory 与 steady outputs？
  \item 与 fixed global thermal basis、trajectory NN 和 high-dimensional online Maxwell solve 相比，精度--训练成本--推理成本的 Pareto 前沿如何？
\end{enumerate}

\subsection{冻结结论}
在 Physics Gate 通过后，主架构固定为
\begin{equation}
\boxed{
 g
 \rightarrow
 \begin{cases}
 \Phi(g),M_r(g),K_r(g)\quad\text{deterministic},\\
 \{Z_{\rm field},D_{\rm vol},H_j\}_\theta\quad\text{neural}
 \end{cases}
 \rightarrow\text{exact current/circuit physics}
 \rightarrow\text{true geometry-aware thermal ROM}.
}
\end{equation}

后续可以改进 POD、network capacity、sampling、active learning、verified bounds、positivity-preserving thermal method、integrator，以及在 held-out evidence 支持时加入 deterministic scale-aware thermal transform；但不得重新把 thermal basis 交给黑箱网络，也不得退回用固定 global span 记忆所有移动 hotspot 的路线。
